% 用 xelatex 编译
%\documentclass[a4paper,cs4size,hyperref,fancyhdr]{ctexbook}
\documentclass[UTF8,openright]{ctexbook}

% 论文版面要求:
% 统一按 word 格式A4纸（页面设置按word默认值）编排、打印、制作。
% 正文内容字体为宋体；字号为小4号；字符间距为标准；行距为25磅(约0.88175cm)。

%%%%% ===== 页面设置
\usepackage[a4paper,top=2.54cm,bottom=2.54cm,left=3.17cm,right=3.17cm,%
            includehead,includefoot]{geometry}
\setlength{\parindent}{2em}
\setlength{\parskip}{0.5em}

%%%%% =====章节 标题 设置
\ctexset{%
  contentsname={目\ \ 录},
  listfigurename={插\ 图\ 目\ 录},
  listtablename={表\ 格\ 目\ 录},
 % bibname={\centerline{参\ 考\ 文\ 献}},
  chapter={name={第,章},
           number=\chinese{chapter},
           nameformat={\zihao{3}\bfseries},
           titleformat={\zihao{3}\bfseries},
           beforeskip={-10pt},
           afterskip={20pt}
           },
  section={format=\raggedright,
           nameformat={\zihao{4}\bfseries},
           titleformat={\zihao{4}\bfseries},
           afterskip={1ex plus 0.2ex}
           },
  subsection={format=\raggedright,
           nameformat={\zihao{-4}\bfseries},
           titleformat={\zihao{-4}\bfseries},
           afterskip={0.5ex plus 0.1ex}
           }
}
%%%%% ===== 中英文字体
\setmainfont{Times New Roman}
%\setsansfont{Myriad Pro} % 无衬线字体 sans serif \sffamily
%\setmonofont{Consolas}   % 等宽字体 typewriter \ttfamily
%\newcommand{\Times}{\fontspec{Times New Roman}}
%% 中文字体
%\setCJKmainfont[BoldFont={Microsoft YaHei},ItalicFont={KaiTi}]{NSimSun}
%\setCJKsansfont{Microsoft YaHei}
%\setCJKmonofont{KaiTi}
%\setCJKfamilyfont{STSong}{方正小标宋_GBK}\newcommand{\STSong}{\CJKfamily{STSong}}
\setCJKfamilyfont{STSong}{STZhongsong}\newcommand{\STSong}{\CJKfamily{STSong}}

%%%%% ===== 常用宏包
\usepackage{amsmath,amssymb,amsfonts,bm}
\usepackage[amsmath,thref,thmmarks,hyperref]{ntheorem}
\usepackage{graphicx,xcolor,float}
\usepackage{fancyhdr}
\usepackage{tocloft} % 设置目录中的条目间距
\usepackage{booktabs} % toprule, midrule, bottomrule
\usepackage{varwidth} % 可变宽度的 parbox

%%%%% ===== 参考文献与链接
\usepackage[numbers,sort&compress,sectionbib]{natbib}
\usepackage[xetex,pagebackref]{hyperref}
\hypersetup{CJKbookmarks=true,colorlinks=true,citecolor=blue,%
            linkcolor=blue,urlcolor=blue,bookmarksnumbered=true,%
	        bookmarksopen=true,breaklinks=true}
\renewcommand*{\backrefalt}[4]{%
\ifcase #1 No citations.%
\or Cited on page #2.%
\else Cited on pages #2.%
%\else #1 Cited on pages #2.%
\fi
}

%%%%% ===== 浮动图表的标题
\usepackage[margin=2em,labelsep=space,skip=0.5em,font=normalfont]{caption}
\DeclareCaptionFormat{mycaption}{{\heiti\color{blue} #1}#2{\kaishu #3}}
\captionsetup{format=mycaption,tablewithin=chapter,figurewithin=chapter}

%%%%% ===== 算法
\usepackage{algorithm,algpseudocode}

%%%%% ===== 其他
\usepackage{ulem}
\def\ULthickness{1pt}

%%%%% ===== 自定义命令
\renewcommand{\C}{\mathbb{C}}
\newcommand{\R}{\mathbb{R}}


%%%%%%%%%%%%%%%%%%%%%%%%%%%%%%%%%%%%%%%%%%%%%%%%%%%%%%%%%%%%%%%%%%%%%
\begin{document}

%%%%% ===== 格式
% !Mode:: "TeX:UTF-8"
%  Authors: 张井   Jing Zhang: prayever@gmail.com     天津大学2010级管理与经济学部信息管理与信息系统专业硕士生
%           余蓝涛 Lantao Yu: lantaoyu1991@gmail.com  天津大学2008级精密仪器与光电子工程学院测控技术与仪器专业本科生

%%%%%%%%%%%%%%%%% Fonts Definition and Basics %%%%%%%%%%%%%%%%%
\setCJKfamilyfont{思源黑体}{Noto Sans CJK SC}
\setCJKfamilyfont{思源宋体}{Noto Serif CJK SC}
\newcommand{\hei}{\CJKfamily{思源黑体}}
\newcommand{\song}{\CJKfamily{思源宋体}}
\newcommand{\yihao}{\fontsize{26pt}{26pt}\selectfont}       % 一号, 单倍行距
\newcommand{\xiaoyi}{\fontsize{24pt}{24pt}\selectfont}      % 小一, 单倍行距
\newcommand{\erhao}{\fontsize{22pt}{1.25\baselineskip}\selectfont}       % 二号, 1.25倍行距
\newcommand{\xiaoer}{\fontsize{18pt}{18pt}\selectfont}      % 小二, 单倍行距
\newcommand{\sanhao}{\fontsize{16pt}{16pt}\selectfont}      % 三号, 单倍行距
\newcommand{\xiaosan}{\fontsize{15pt}{15pt}\selectfont}     % 小三, 单倍行距
\newcommand{\sihao}{\fontsize{14pt}{14pt}\selectfont}       % 四号, 单倍行距
\newcommand{\xiaosi}{\fontsize{12pt}{12pt}\selectfont}      % 小四, 单倍行距
\newcommand{\wuhao}{\fontsize{10.5pt}{10.5pt}\selectfont}   % 五号, 单倍行距
\newcommand{\xiaowu}{\fontsize{9pt}{9pt}\selectfont}        % 小五, 单倍行距

%\CJKtilde  % 重新定义了波浪符~的意义
% JUST DON'T USE CJK
% 使用 ctexbook 之后已无必要
\newcommand\prechaptername{第}
\newcommand\postchaptername{章}

\punctstyle{hangmobanjiao}             % 调整中文字符的表示，行内占一个字符宽度，行尾占半个字符宽度

% 调整罗列环境的布局
\setitemize{leftmargin=3em,itemsep=0em,partopsep=0em,parsep=0em,topsep=-0em}
\setenumerate{leftmargin=3em,itemsep=0em,partopsep=0em,parsep=0em,topsep=0em}

% 避免宏包 hyperref 和 arydshln 不兼容带来的目录链接失效的问题。
\def\temp{\relax}
\let\temp\addcontentsline
\gdef\addcontentsline{\phantomsection\temp}

% 自定义项目列表标签及格式 \begin{publist} 列表项 \end{publist}
\newcounter{pubctr} %自定义新计数器
\newenvironment{publist}{%%%%%定义新环境
\begin{list}{[\arabic{pubctr}]} %%标签格式
    {
     \usecounter{pubctr}
     \setlength{\leftmargin}{2.5em}   % 左边界 \leftmargin =\itemindent + \labelwidth + \labelsep
     \setlength{\itemindent}{0em}     % 标号缩进量
     \setlength{\labelsep}{1em}       % 标号和列表项之间的距离,默认0.5em
     \setlength{\rightmargin}{0em}    % 右边界
     \setlength{\topsep}{0ex}         % 列表到上下文的垂直距离
     \setlength{\parsep}{0ex}         % 段落间距
     \setlength{\itemsep}{0ex}        % 标签间距
     \setlength{\listparindent}{0pt}  % 段落缩进量
    }}
{\end{list}}

\makeatletter
\renewcommand\normalsize{
  \@setfontsize\normalsize{12pt}{12pt} % 小四对应 12 pt
  \setlength\abovedisplayskip{4pt}
  \setlength\abovedisplayshortskip{4pt}
  \setlength\belowdisplayskip{\abovedisplayskip}
  \setlength\belowdisplayshortskip{\abovedisplayshortskip}
  \let\@listi\@listI}
\def\defaultfont{\renewcommand{\baselinestretch}{1.63}\normalsize\selectfont} % 设置行距

\renewcommand{\CJKglue}{\hskip -0.1 pt plus 0.08\baselineskip} % 控制字间距，使每行 34 个汉字
\makeatother

%%%%%%%%%%%%% Contents %%%%%%%%%%%%%%%%%
\renewcommand{\contentsname}{目\qquad 录}
\setcounter{tocdepth}{2} % 控制目录深度
% 使用 ctexbook 之后已无必要
%\titlecontents{chapter}[2em]{\vspace{.5\baselineskip}\xiaosan\song}
             %{\prechaptername\CJKnumber{\thecontentslabel}\postchaptername\qquad}{}
             %{\hspace{.5em}\titlerule*[10pt]{$\cdot$}\sihao\contentspage}
\titlecontents{chapter}[2em]{\vspace{.5\baselineskip}\xiaosan\song}
             {\thecontentslabel\!\!\qquad}{}
             {\titlerule*[5pt]{$\cdot$}\sihao\contentspage}
\titlecontents{section}[3em]{\vspace{.25\baselineskip}\sihao\song}
             {\thecontentslabel\quad}{}
             {\titlerule*[5pt]{$\cdot$}\sihao\contentspage}
\titlecontents{subsection}[4em]{\vspace{.25\baselineskip}\sihao\song}
             {\thecontentslabel\quad}{}
             {\titlerule*[5pt]{$\cdot$}\sihao\contentspage}

%%%%%%%%%% Chapter and Section %%%%%%%%%%%%%
\setcounter{secnumdepth}{4}
\setlength{\parindent}{2em}
\renewcommand{\chaptername}{\prechaptername\CJKnumber{\thechapter}\postchaptername}
\titleformat{\chapter}{\centering\xiaosan\hei}{\chaptername}{2em}{}
\titlespacing{\chapter}{0pt}{0.1\baselineskip}{0.8\baselineskip}
\titleformat{\section}{\sihao\hei}{\thesection}{1em}{}
\titlespacing{\section}{0pt}{0.15\baselineskip}{0.25\baselineskip}
\titleformat{\subsection}{\sihao\hei}{\thesubsection}{1em}{}
\titlespacing{\subsection}{0pt}{0.1\baselineskip}{0.3\baselineskip}
\titleformat{\subsubsection}{\sihao\hei}{\thesubsubsection}{1em}{}
\titlespacing{\subsubsection}{0pt}{0.05\baselineskip}{0.1\baselineskip}

\newcommand{\trtitle}[1]{\vspace{12pt}\begin{center}{\xiaoer\hei\bf #1}\end{center}}
\newcommand{\trchapter}[1]{\vspace{6pt}{\raggedright\xiaosan\hei #1}\newline}
\newcommand{\trsection}[1]{\vspace{4pt}{\raggedright\sihao\hei #1}\newline}
\newcommand{\trsubsection}[1]{\vspace{2pt}{\raggedright\sihao\hei\it #1}\newline}

%%%%%%%%%% Table, Figure and Equation %%%%%%%%%%%%%%%%%
\renewcommand{\tablename}{表}                                     % 插表题头
\renewcommand{\figurename}{图}                                    % 插图题头
\renewcommand{\thefigure}{\arabic{chapter}-\arabic{figure}}       % 使图编号为 7-1 的格式 %\protect{~}
\renewcommand{\thesubfigure}{\alph{subfigure})}                   % 使子图编号为 a) 的格式
\renewcommand{\thesubtable}{(\alph{subtable})}                    % 使子表编号为 (a) 的格式
\renewcommand{\thetable}{\arabic{chapter}-\arabic{table}}         % 使表编号为 7-1 的格式
\renewcommand{\theequation}{\arabic{chapter}-\arabic{equation}}   % 使公式编号为 7-1 的格式
\newcommand{\ud}{\mathrm{d}}

%%%%%% 定制浮动图形和表格标题样式 %%%%%%
\makeatletter
\long\def\@makecaption#1#2{
   \vskip\abovecaptionskip
   \sbox\@tempboxa{\centering\wuhao\song{#1\qquad #2} }
   \ifdim \wd\@tempboxa >\hsize
     \centering\wuhao\song{#1\qquad #2} \par
   \else
     \global \@minipagefalse
     \hb@xt@\hsize{\hfil\box\@tempboxa\hfil}
   \fi
   \vskip\belowcaptionskip\vspace{8pt}}
\makeatother
\captiondelim{~~~~} %用来控制longtable表头分隔符

%%%%%%%%%% Theorem Environment %%%%%%%%%%%%%%%%%
\theoremstyle{plain}
\theorembodyfont{\song\rmfamily}
\theoremheaderfont{\hei\rmfamily}
\newtheorem{theorem}{定理~}[chapter]
\newtheorem{lemma}{引理~}[chapter]
\newtheorem{axiom}{公理~}[chapter]
\newtheorem{proposition}{命题~}[chapter]
\newtheorem{prop}{性质~}[chapter]
\newtheorem{corollary}{推论~}[chapter]
\newtheorem{definition}{定义~}[chapter]
\newtheorem{conjecture}{猜想~}[chapter]
\newtheorem{example}{例~}[chapter]
\newtheorem{remark}{注~}[chapter]
%\newtheorem{algorithm}{算法~}[chapter]
\newenvironment{proof}{\noindent{\hei 证明：}}{\hfill $ \square $ \vskip 4mm}
\theoremsymbol{$\square$}

%%%%%%%%%% Page: number, header and footer  %%%%%%%%%%%%%%%%%

%\frontmatter 或 \pagenumbering{roman}
%\mainmatter 或 \pagenumbering{arabic}
\makeatletter
\renewcommand\frontmatter{\clearpage
  \@mainmatterfalse
  }
\makeatother

%%%%%%%%%%% Code: Listings from MCM Template %%%%%%%%%%%%

\definecolor{grey}{rgb}{0.8,0.8,0.8}
\definecolor{darkgreen}{rgb}{0,0.3,0}
\definecolor{darkblue}{rgb}{0,0,0.3}
\def\lstbasicfont{\fontfamily{pcr}\selectfont\footnotesize}
\lstset{%
% indexing
   % numbers=left,
   % numberstyle=\small,%
% character display
    showstringspaces=false,
    showspaces=false,%
    tabsize=4,%
% style
    frame=lines,%
    basicstyle={\footnotesize\lstbasicfont},%
    keywordstyle=\color{darkblue}\bfseries,%
    identifierstyle=,%
    commentstyle=\color{darkgreen},%\itshape,%
    stringstyle=\color{black}%
}
\lstloadlanguages{C,C++,Java,Matlab,Mathematica,Python}

%%%%%%%%%%%% References %%%%%%%%%%%%%%%%%

\makeatletter
\renewenvironment{thebibliography}[1]{
   \wuhao
   \list{\@biblabel{\@arabic\c@enumiv}}
        {\renewcommand{\makelabel}[1]{##1\hfill}
         \settowidth\labelwidth{0 cm}
         \setlength{\labelsep}{0pt}
         \setlength{\itemindent}{0pt}
         \setlength{\leftmargin}{\labelwidth+\labelsep}
         \addtolength{\itemsep}{-0.7em}
         \usecounter{enumiv}
         \let\p@enumiv\@empty
         \renewcommand\theenumiv{\@arabic\c@enumiv}}
    \sloppy\frenchspacing
    \clubpenalty4000
    \@clubpenalty \clubpenalty
    \widowpenalty4000
    \interlinepenalty4000
    \sfcode`\.\@m}
   {\def\@noitemerr
     {\@latex@warning{Empty 'thebibliography' environment}}
    \endlist\frenchspacing}
\makeatother

\addtolength{\bibsep}{-0.5em}     % 缩小参考文献间的垂直间距
\setlength{\bibhang}{2em}         % 每个条目自第二行起缩进的距离

% 参考文献引用作为上标出现
\newcommand{\citeup}[1]{\textsuperscript{\cite{#1}}}

%%%%%%%%%%%% Cover %%%%%%%%%%%%%%%%%
% 封面、摘要、版权、致谢格式定义
\makeatletter
\def\ctitle#1{\def\@ctitle{#1}}\def\@ctitle{}
\def\cdegree#1{\def\@cdegree{#1}}\def\@cdegree{}
\def\caffil#1{\def\@caffil{#1}}\def\@caffil{}
\def\csubject#1{\def\@csubject{#1}}\def\@csubject{}
\def\cgrade#1{\def\@cgrade{#1}}\def\@cgrade{}
\def\cauthor#1{\def\@cauthor{#1}}\def\@cauthor{}
\def\cnumber#1{\def\@cnumber{#1}}\def\@cnumber{}
\def\csupervisor#1{\def\@csupervisor{#1}}\def\@csupervisor{}
\def\crank#1{\def\@crank{#1}}\def\@crank{}
\def\cdate#1{\def\@cdate{#1}}\def\@cdate{}
\long\def\cabstract#1{\long\def\@cabstract{#1}}\long\def\@cabstract{}
\long\def\eabstract#1{\long\def\@eabstract{#1}}\long\def\@eabstract{}
\def\ckeywords#1{\def\@ckeywords{#1}}\def\@ckeywords{}
\def\ekeywords#1{\def\@ekeywords{#1}}\def\@ekeywords{}
\def\cheading#1{\def\@cheading{#1}}\def\@cheading{}


\pagestyle{fancy}
\fancyhf{}
\fancyhead[C]{\song\wuhao \@cheading}  % 页眉显示天津大学 20XX 届本科生毕业论文
\fancyfoot[C]{\song\xiaowu ~\thepage~}
\newlength{\@title@width}

%%%%%%%%%%%%%%%%%%%     Cover     %%%%%%%%%%%%%%%%%%%%%%%
\def\makecover{
   	\phantomsection
    \pdfbookmark[-1]{\@ctitle}{ctitle}

    \begin{titlepage}
      	\vspace*{31.5pt}
      	\begin{center}
  			\begin{figure}[h]
  				\centering
  				\includegraphics[width=0.4\textwidth]{figures/tjuname}
  			\end{figure}
  			\vspace*{21pt}
  			\hei\erhao{\textbf{本科生毕业论文}}
  			\vspace*{52.5pt}

  			\begin{figure}[h]
  				\centering
  				\includegraphics[width=0.4\textwidth]{figures/tjulogo}
  			\end{figure}
  			\vspace*{42pt}
  			
  			\renewcommand\arraystretch{1.5}
  			\setlength{\@title@width}{5cm}{\song\sanhao{\textbf{
    			\begin{tabular}{lc}
      				学\qquad 院&	\underline{\makebox[\@title@width][c]{\@caffil}}	\\
      				专\qquad 业&	\underline{\makebox[\@title@width][c]{\@csubject}}	\\
      				年\qquad 级&	\underline{\makebox[\@title@width][c]{\@cgrade}}	\\
      				姓\qquad 名&	\underline{\makebox[\@title@width][c]{\@cauthor}} 	\\
      				指导教师   &	\underline{\makebox[\@title@width][c]{\@csupervisor}} \\
  				\end{tabular}
  			}}}
  			\vspace*{21pt}

			\song\sanhao{\textbf{\@cdate}}
		\end{center}
	\end{titlepage}
}

%%%%%%%%%%%%%%%%%%%    Abstract    %%%%%%%%%%%%%%%%%%%%%%%
\def\makeabstract{
	\titleformat{\chapter}{\centering\sanhao\hei}{\chaptername}{2em}{} %"摘要"三号黑体居中

%%%%% Chinese Abstract %%%%%
	\clearpage
	\markboth{摘~要}{摘~要}
	\pdfbookmark[0]{摘~~要}{cabstract}
	\chapter*{摘\qquad 要}
	
	\song\defaultfont
	\@cabstract
	\vspace{\baselineskip}

	\hangafter=1\hangindent=52.3pt\noindent
	{\hei\xiaosi 关键词：} \@ckeywords
	\thispagestyle{empty}

%%%%% English Abstract %%%%%
	\clearpage
	\markboth{ABSTRACT}{ABSTRACT}
	\pdfbookmark[0]{ABSTRACT}{eabstract}
	\chapter*{\bf{ABSTRACT}}
	
	\@eabstract
	\vspace{\baselineskip}

	\hangafter=1\hangindent=60pt\noindent
	{\textbf{Keywords:}} \@ekeywords
	\thispagestyle{empty}
}

%%%%%%%%%%%%%%%%%%%    Contents    %%%%%%%%%%%%%%%%%%%%%%%
\def\makecontents{
	\titleformat{\chapter}{\centering\sanhao\hei}{\chaptername}{2em}{} %"目录"三号黑体居中

	\clearpage{\pagestyle{empty}\cleardoublepage}
	\pdfbookmark[0]{目~~录}{mulu}
	\tableofcontents                                     % 中文目录
	\fancypagestyle{plain}{
		\fancyhf{}
		\renewcommand{\headrulewidth}{0 pt}
		\fancyfoot[C]{\song\xiaowu~\thepage~}
	}
	\thispagestyle{plain}

	\mainmatter\defaultfont\sloppy\raggedbottom
	\makeatletter
	\fancypagestyle{plain}{                              % 设置开章页眉页脚风格
    	\fancyhf{}
    	\fancyhead[C]{\song\wuhao \@cheading}            % 首页页眉格式
    	\fancyfoot[C]{\song\xiaowu ~\thepage~}           % 首页页脚格式
    	\renewcommand{\headrulewidth}{0.5pt}
    	\renewcommand{\footrulewidth}{0pt}
	}
}

%%%%%%%%%%%%%%%%%%%   References   %%%%%%%%%%%%%%%%%%%%%%%
\def\makereferences{
	\titleformat{\chapter}{\raggedright\sihao\hei}{\chaptername}{2em}{} % "参考文献"四号黑体顶格
	\chapter*{\bibname}
	\defaultfont
	\bibliographystyle{references/ref.bst}
	\bibliography{references/ref}

	\phantomsection
	\markboth{参考文献}{参考文献}
	\addcontentsline{toc}{chapter}{参考文献}       % 参考文献加入到中文目录
	\nocite{*}                                     % 若将此命令屏蔽掉，则未引用的文献不会出现在文后的参考文献中
}

\makeatother


\zihao{-4}

%%%%% ===== 中文封面信息
\graduateyear{2016} % 毕业年份
\class{O241.6} % 分类号（数值线性代数是 O241.6）
\ctitle{\uline{论文标题论文标题标题\linebreak 如果一行放不下就放两行}}
\def\cctitle{论文标题论文标题} % 在原创性声明中使用, 不能出现手工换行
\caffil{数学系}
\cmajor{XXXX} % 计算数学
\cdirection{XXXX} % 数值代数
\csupervisor{某某某\, 教授}
\cauthor{XXXX}
\def\ccauthor{\@cauthor} % 在答辩委员会中使用
\studentid{51088888888}
\cdate{20xx 年 xx 月}

%%%%% ===== 英文封面信息
\etitle{\uline{Title of Thesis Title of Thesis\linebreak
       Title Title Title of Thesis Title of\linebreak
       Thesis Title Title Title Title Title}}
\eaffil{Mathematics} 
\emajor{xxxx xxxx} % Computational Mathematics
\edirection{xxxx xxxx} % Numerical Algebra
\esupervisor{XXX Xxxxxxx (Professor)}
\eauthor{ZHANG San}
\edate{xxx, 20xx}

%%%%% ===== 生成封面 =====
\newgeometry{top=2.0cm,bottom=2.0cm,left=2.5cm,right=2.5cm}
{
\renewcommand{\baselinestretch}{1.6}
\makecover
}

%%%%% ===== 原创性声明与著作权使用声明 =====
%%%%% ===== 原创性声明与著作权使用声明
\newpage
\thispagestyle{empty}
\pdfbookmark[0]{原创性声明}{Declaration}

\vspace*{1em}

{
\linespread{1.4}\zihao{-4}
\centerline{\zihao{3}\STSong 华东师范大学学位论文原创性声明}

\bigskip

郑重声明：本人呈交的学位论文《\cctitle》，
是在华东师范大学攻读硕士/博士（请勾选）
学位期间，在导师的指导下进行的研究工作及取得的研究成果。
除文中已经注明引用的内容外，本论文不包含其他个人已经发表或撰写过的研究成果。
对本文的研究做出重要贡献的个人和集体，均已在文中作了明确说明并表示谢意。

\vspace{1em}


{\STSong 作者签名}：$\underline{\hspace{4cm}}$ \hfill
{\STSong 日\quad 期}： \qquad\ 年 \quad\ 月 \quad\ 日 \qquad\mbox{}

\vspace{4em}

\centerline{\zihao{3}\STSong 华东师范大学学位论文著作权使用声明}
\bigskip


《\cctitle》
系本人在华东师范大学攻读学位期间在导师指导下完成的硕士/博士（请勾选）学位论文，
本论文的著作权归本人所有。
本人同意华东师范大学根据相关规定保留和使用此学位论文，
并向主管部门和学校指定的相关机构送交学位论文的印刷版和电子版；
允许学位论文进入华东师范大学图书馆及数据库被查阅、借阅；
同意学校将学位论文加入全国博士、硕士学位论文共建单位数据库进行检索，
将学位论文的标题和摘要汇编出版，采用影印、缩印或者其它方式合理复制学位论文。

本学位论文属于（请勾选）
\newcounter{muni}
\begin{list}{{\hfill\upshape (\qquad)\ \arabic{muni}. }}{%
     \usecounter{muni}\leftmargin6.5em\labelwidth4.2em\labelsep0.2em
     \itemsep0.5em\itemindent0pt\parsep0pt\topsep0pt}

\item 经华东师范大学相关部门审查核定的“内部”或“涉密”学位论文*，
  于 \qquad 年 \quad 月 \quad  日解密, 解密后适用上述授权。

\item 不保密，适用上述授权。
\end{list}

\vskip0.8cm


{\STSong 导师签名}：$\underline{\hspace{4cm}}$ \hfill
{\STSong 本人签名}：$\underline{\hspace{4cm}}$

\bigskip

{\mbox{}\hfill 年\qquad 月\qquad  日 }

\vfill

\parbox[t]{0.946\textwidth}{\zihao{5}
*“涉密”学位论文应是已经华东师范大学学位评定委员会办公室或保密委员会
审定过的学位论文（需附获批的《华东师范大学研究生申请学位论文“涉密”审批表》
方为有效），未经上述部门审定的学位论文均为公开学位论文。
此声明栏不填写的，默认为公开学位论文，均适用上述授权）。\\
}}


%%%%% ===== 答辩委员会成员 =====
%%%% ===== 答辩委员会成员
\thispagestyle{empty}
\pdfbookmark[0]{答辩委员会}{Committee}
\vspace*{2em}

\makeatletter
\begin{center}\STSong\zihao{3}
 \underline{\ \@cauthor\ } 硕士学位论文答辩委员会成员名单
\end{center}
\makeatother

\begin{center}\zihao{4}
\renewcommand{\arraystretch}{1.4}
  \begin{tabular}{|c|c|c|c|} \hline
   ~~~~~姓~名~~~~~ & ~~~~~职~称~~~~~
   & \hspace{6em}单~位\hspace{6em} & ~~~备~注~~~\\\hline
         XXX    &   教授    &  XXXXX大学数学系  & 主席  \\ \hline
         XXX    &   教授    &  XXXXX大学数学系  &       \\ \hline
         XXX    &   教授    &  XXXXX大学数学系  &       \\ \hline
                &           &                   &       \\ \hline
                &           &                   &       \\ \hline
  \end{tabular}
\end{center}


\frontmatter
%\restoregeometry
%%%%% ===== 中文摘要 =====
\include{Abstract_chs}

%%%%% ===== 英文摘要 =====
\include{Abstract_eng}

%%%%% ===== 生成目录
\setcounter{tocdepth}{1}
\phantomsection
\addcontentsline{toc}{chapter}{目录}
\tableofcontents
%\listoffigures % 插图目录
%\listoftables  % 表格目录
\clearpage%{\pagestyle{empty}\cleardoublepage}

%%%%%% ===== 正文部分 ===== %%%%%
\mainmatter
\linespread{1.4}\selectfont
%\setlength{\baselineskip}{0.88175cm}

\chapter{引言}

引言部分, 介绍论文研究课题的应用背景或者问题来源,
一些基本概念, 现有成果等等.

\section{问题的提出}
问题的提出, 问题的提出, 问题的提出, 问题的提出, 问题的提出.
问题的提出, 问题的提出, 问题的提出, 问题的提出, 问题的提出.
问题的提出, 问题的提出, 问题的提出, 问题的提出, 问题的提出.

问题的提出, 问题的提出, 问题的提出, 问题的提出, 问题的提出.
问题的提出, 问题的提出, 问题的提出, 问题的提出, 问题的提出.
问题的提出, 问题的提出, 问题的提出, 问题的提出, 问题的提出.

\section{现有成果}
现有成果介绍, 现有成果介绍, 现有成果介绍, 现有成果介绍, 现有成果介绍.
现有成果介绍, 现有成果介绍, 现有成果介绍, 现有成果介绍, 现有成果介绍.
现有成果介绍, 现有成果介绍, 现有成果介绍, 现有成果介绍, 现有成果介绍.

\begin{theorem}
  这是定理, 这是定理, 这是定理, 这是定理, 这是定理, 这是定理.
  这是定理, 这是定理, 这是定理, 这是定理, 这是定理, 这是定理.
  这是定理, 这是定理, 这是定理, 这是定理, 这是定理, 这是定理.
  这是定理, 这是定理, 这是定理, 这是定理, 这是定理, 这是定理.
\end{theorem}
\begin{proof}
  这是证明, 这是证明, 这是证明, 这是证明, 这是证明, 这是证明.
  这是证明, 这是证明, 这是证明, 这是证明, 这是证明, 这是证明.
  这是证明, 这是证明, 这是证明, 这是证明, 这是证明, 这是证明.
  这是证明, 这是证明, 这是证明, 这是证明, 这是证明, 这是证明.

  这是证明, 这是证明, 这是证明, 这是证明, 这是证明, 这是证明.
  这是证明, 这是证明, 这是证明, 这是证明, 这是证明, 这是证明.
  这是证明, 这是证明, 这是证明, 这是证明, 这是证明, 这是证明.
  这是证明, 这是证明, 这是证明, 这是证明, 这是证明, 这是证明.
\end{proof}

现有成果介绍, 现有成果介绍, 现有成果介绍, 现有成果介绍, 现有成果介绍.
现有成果介绍, 现有成果介绍, 现有成果介绍, 现有成果介绍, 现有成果介绍.
现有成果介绍, 现有成果介绍, 现有成果介绍, 现有成果介绍, 现有成果介绍.
现有成果介绍, 现有成果介绍, 现有成果介绍, 现有成果介绍, 现有成果介绍.
现有成果介绍, 现有成果介绍, 现有成果介绍, 现有成果介绍, 现有成果介绍.
现有成果介绍, 现有成果介绍, 现有成果介绍, 现有成果介绍, 现有成果介绍.

\chapter{准备工作}

论文中需要用到的一些知识或工具等. 论文中需要用到的一些知识或工具等.
论文中需要用到的一些知识或工具等. 论文中需要用到的一些知识或工具等.
论文中需要用到的一些知识或工具等. 论文中需要用到的一些知识或工具等.
论文中需要用到的一些知识或工具等. 论文中需要用到的一些知识或工具等.
论文中需要用到的一些知识或工具等. 论文中需要用到的一些知识或工具等.

\section{一些定义}
这里给出一些定义, 这里给出一些定义, 这里给出一些定义,
这里给出一些定义, 这里给出一些定义, 这里给出一些定义.

\begin{definition}
  这是定义, 这是定义, 这是定义, 这是定义, 这是定义, 这是定义,
  这是定义, 这是定义, 这是定义, 这是定义, 这是定义, 这是定义,
  这是定义, 这是定义, 这是定义, 这是定义, 这是定义, 这是定义,
  这是定义, 这是定义, 这是定义, 这是定义, 这是定义, 这是定义.
\end{definition}

\section{一些引理和推论}

这里给出一些引理和推论, 这里给出一些引理和推论, 这里给出一些引理和推论,
这里给出一些引理和推论, 这里给出一些引理和推论, 这里给出一些引理和推论,
这里给出一些引理和推论, 这里给出一些引理和推论, 这里给出一些引理和推论.

\begin{lemma}
  这是引理, 这是引理, 这是引理, 这是引理, 这是引理, 这是引理,
  这是引理, 这是引理, 这是引理, 这是引理, 这是引理, 这是引理,
  这是引理, 这是引理, 这是引理, 这是引理, 这是引理, 这是引理,
  这是引理, 这是引理, 这是引理, 这是引理, 这是引理, 这是引理.
\end{lemma}

\begin{lemma}
  这是引理, 这是引理, 这是引理, 这是引理, 这是引理, 这是引理,
  这是引理, 这是引理, 这是引理, 这是引理, 这是引理, 这是引理,
  这是引理, 这是引理, 这是引理, 这是引理, 这是引理, 这是引理,
  这是引理, 这是引理, 这是引理, 这是引理, 这是引理, 这是引理.
\end{lemma}

\subsection{推论}
\begin{corollary}
  这是推论, 这是推论, 这是推论, 这是推论, 这是推论, 这是推论,
  这是推论, 这是推论, 这是推论, 这是推论, 这是推论, 这是推论,
  这是推论, 这是推论, 这是推论, 这是推论, 这是推论, 这是推论,
  这是推论, 这是推论, 这是推论, 这是推论, 这是推论, 这是推论.
\end{corollary}

\section{一些命题和性质}
现有成果介绍, 现有成果介绍, 现有成果介绍, 现有成果介绍, 现有成果介绍.
现有成果介绍, 现有成果介绍, 现有成果介绍, 现有成果介绍, 现有成果介绍.
现有成果介绍, 现有成果介绍, 现有成果介绍, 现有成果介绍, 现有成果介绍.

\begin{proposition}
  这是命题, 这是命题, 这是命题, 这是命题, 这是命题, 这是命题,
  这是命题, 这是命题, 这是命题, 这是命题, 这是命题, 这是命题,
  这是命题, 这是命题, 这是命题, 这是命题, 这是命题, 这是命题.
\end{proposition}

\begin{property}
  这是性质, 这是性质, 这是性质, 这是性质, 这是性质, 这是性质,
  这是性质, 这是性质, 这是性质, 这是性质, 这是性质, 这是性质,
  这是性质, 这是性质, 这是性质, 这是性质, 这是性质, 这是性质.
\end{property}

现有成果介绍, 现有成果介绍, 现有成果介绍, 现有成果介绍, 现有成果介绍.
现有成果介绍, 现有成果介绍, 现有成果介绍, 现有成果介绍, 现有成果介绍.
现有成果介绍, 现有成果介绍, 现有成果介绍, 现有成果介绍, 现有成果介绍.


\section{一些注记}
现有成果介绍, 现有成果介绍, 现有成果介绍, 现有成果介绍, 现有成果介绍.
现有成果介绍, 现有成果介绍, 现有成果介绍, 现有成果介绍, 现有成果介绍.
现有成果介绍, 现有成果介绍, 现有成果介绍, 现有成果介绍, 现有成果介绍.
现有成果介绍, 现有成果介绍, 现有成果介绍, 现有成果介绍, 现有成果介绍.
现有成果介绍, 现有成果介绍, 现有成果介绍, 现有成果介绍, 现有成果介绍.

\begin{remark}
  这是注记, 这是注记, 这是注记, 这是注记, 这是注记, 这是注记,
  这是注记, 这是注记, 这是注记, 这是注记, 这是注记, 这是注记,
  这是注记, 这是注记, 这是注记, 这是注记, 这是注记, 这是注记.
\end{remark}

\section{算法}
算法示例, 算法示例, 算法示例, 算法示例, 算法示例, 算法示例,
算法示例, 算法示例, 算法示例, 算法示例, 算法示例, 算法示例,
算法示例, 算法示例, 算法示例, 算法示例, 算法示例, 算法示例.

算法示例, 算法示例, 算法示例, 算法示例, 算法示例, 算法示例,
算法示例, 算法示例, 算法示例, 算法示例, 算法示例, 算法示例,
算法示例, 算法示例, 算法示例, 算法示例, 算法示例, 算法示例.

算法示例, 算法示例, 算法示例, 算法示例, 算法示例, 算法示例,
算法示例, 算法示例, 算法示例, 算法示例, 算法示例, 算法示例,
算法示例, 算法示例, 算法示例, 算法示例, 算法示例, 算法示例.

\begin{algorithm}[H]
\caption{算法示例算法示例\label{Alg:}}
  \begin{algorithmic}[1]
  \State $v_1 = r/\|r\|_2$
  \For{$j=1$ to $m$}
  \State $z = A v_j$
  \For{$i=1$ to $j$}
  \State $h_{i,j} = (v_i,z)$  \Comment{内积}
  \State $z = z - h_{i,j}v_i$
  \EndFor
  \State $h_{j+1,j}=\|z\|_2$
  \If{$h_{j+1,j}=0$}
  \State break
  \EndIf
  \State $v_{j+1}=z/h_{j+1,j}$
  \EndFor
  \end{algorithmic}
\end{algorithm}

\clearpage{\pagestyle{empty}\cleardoublepage}
\chapter{其他章节}
其他章节, 其他章节, 其他章节, 其他章节, 其他章节, 其他章节,
其他章节, 其他章节, 其他章节, 其他章节, 其他章节, 其他章节,
其他章节, 其他章节, 其他章节, 其他章节, 其他章节, 其他章节.

其他章节, 其他章节, 其他章节, 其他章节, 其他章节, 其他章节,
其他章节, 其他章节, 其他章节, 其他章节, 其他章节, 其他章节,
其他章节, 其他章节, 其他章节, 其他章节, 其他章节, 其他章节.


\clearpage{\pagestyle{empty}\cleardoublepage}
\backmatter
\linespread{1.1}\selectfont
%%%% ===== 参考文献 =====
\setlength{\bibsep}{1ex}  % 需 natbib 宏包
\begin{thebibliography}{99}
\addcontentsline{toc}{chapter}{参考文献}
\thispagestyle{plain}

% \addtolength{\itemsep}{-5pt}

\bibitem{CPZ08}
  K. C. Chang, K. Pearson and T. Zhang,
  Perron-Frobenius Theorem for nonnegative tensors,
  Commun. Math. Sci., 6 (2008), 507--520.

\bibitem{NW99}
  J. Nocedal and S. J. Wright,
  Numerical Optimization,
  Springer, New York, 1999.

\end{thebibliography}

\clearpage{\pagestyle{empty}\cleardoublepage}
\linespread{1.4}\selectfont
\chapter*{附录}
\addcontentsline{toc}{chapter}{附录}

附录部分, 附录部分, 附录部分, 附录部分, 附录部分,
附录部分, 附录部分, 附录部分, 附录部分, 附录部分,
附录部分, 附录部分, 附录部分, 附录部分, 附录部分,
附录部分, 附录部分, 附录部分, 附录部分, 附录部分.

附录部分, 附录部分, 附录部分, 附录部分, 附录部分,
附录部分, 附录部分, 附录部分, 附录部分, 附录部分,
附录部分, 附录部分, 附录部分, 附录部分, 附录部分,
附录部分, 附录部分, 附录部分, 附录部分, 附录部分.

附录部分, 附录部分, 附录部分, 附录部分, 附录部分,
附录部分, 附录部分, 附录部分, 附录部分, 附录部分,
附录部分, 附录部分, 附录部分, 附录部分, 附录部分,
附录部分, 附录部分, 附录部分, 附录部分, 附录部分.


\clearpage{\pagestyle{empty}\cleardoublepage}
\chapter*{致谢}
\addcontentsline{toc}{chapter}{致谢}

致谢部分, 致谢部分, 致谢部分, 致谢部分, 致谢部分,
致谢部分, 致谢部分, 致谢部分, 致谢部分, 致谢部分,
致谢部分, 致谢部分, 致谢部分, 致谢部分, 致谢部分,
致谢部分, 致谢部分, 致谢部分, 致谢部分, 致谢部分.

致谢部分, 致谢部分, 致谢部分, 致谢部分, 致谢部分,
致谢部分, 致谢部分, 致谢部分, 致谢部分, 致谢部分,
致谢部分, 致谢部分, 致谢部分, 致谢部分, 致谢部分,
致谢部分, 致谢部分, 致谢部分, 致谢部分, 致谢部分.

致谢部分, 致谢部分, 致谢部分, 致谢部分, 致谢部分,
致谢部分, 致谢部分, 致谢部分, 致谢部分, 致谢部分,
致谢部分, 致谢部分, 致谢部分, 致谢部分, 致谢部分,
致谢部分, 致谢部分, 致谢部分, 致谢部分, 致谢部分.

致谢部分, 致谢部分, 致谢部分, 致谢部分, 致谢部分,
致谢部分, 致谢部分, 致谢部分, 致谢部分, 致谢部分,
致谢部分, 致谢部分, 致谢部分, 致谢部分, 致谢部分,
致谢部分, 致谢部分, 致谢部分, 致谢部分, 致谢部分.

\clearpage{\pagestyle{empty}\cleardoublepage}
\chapter*{研究成果}
\addcontentsline{toc}{chapter}{研究成果}

研究生期间所取得的研究成果.

\end{document}
